\chapter{システムおよび評価方法の改良}

\section{予備評価を踏まえた設計方針}

予備評価では，当初，複数のグリッチ表現を比較することで，多くの利用者に共通して好まれる表現を明らかにすることを目的としていた．しかし，8種類のライブ壁紙に対する評価は特定の一種類へ集中せず，参加者によって高く評価する表現が異なる結果となった．

この結果から，本研究では，単一のグリッチ表現を最適なものとして決定するよりも，利用者が自身の好みに応じて異なる表現を選択できる構成が適している可能性があると考えた．そこで，予備評価以降は，利用者間の選好の違いを考慮したライブ壁紙の設計へ研究方針を変更した．

また，予備評価で使用した8種類のプリセットは，RGBずれ量，ノイズサイズ，パララックス量，ちらつきなど，複数の設定値をそれぞれ異なる組合せで変更したものであった．そのため，各プリセット間の違いが多く，利用者にとって表現の特徴を把握しにくいという問題もあった．

そこで改良版では，表現の違いをより明確にするため，端末操作に対する反応方式と配色という二つの軸を中心としてプリセットを再構成した．端末操作に対する反応方式として，画面の破損や乱れを強調する「ザビ」と，レイヤの移動やパララックスを中心とする「ゆら」を設定し，それぞれにモノクロ系とカラー系の配色を組み合わせた．これにより，「ザビモノ」，「ザビカラー」，「ゆらモノ」，「ゆらカラー」の4種類へ整理した．

なお，この変更は旧8種類のプリセットを単純に統合したものではない．プリセット数の整理と同時に，グリッチ模様の生成方法，端末操作への反応，レイヤ構成，配色処理などについても改良を行った．


\section{グリッチ表現の改良}

予備評価時のシステムにも，矩形ノイズ，帯状ノイズ，RGBずれ，ちらつきなどのグリッチ表現を実装していた．一方，初期版では比較的大きな矩形や規則的な形状が目立ちやすく，複数の図形を重ねたような外観となる場合があった．

そこで改良版では，単純な矩形の配置だけではなく，横長・縦長の細かな破片，輪郭の欠損，内部の穴や切れ目，粒状ノイズ，傷状の模様などを追加した．また，帯状ノイズについても，単純な帯の移動だけでなく，領域ごとの横方向のずれ，上下方向の歪み，帯同士の重なりなどを組み合わせた．

これにより，矩形や線が配置されている印象を弱め，画面の一部が欠ける，裂ける，ずれるといった複数の表示異常が同時に発生しているような外観へ変更した．

また，改良版では「ザビ」と「ゆら」でグリッチの構成そのものを変化させる．

「ザビ」では，細かなノイズを比較的高い密度で配置し，RGBずれ，帯状ノイズ，ちらつき，位置のずれなどを強くする．さらに，前景側に細かな破片を追加し，端末を動かした際に画面上の模様がその場で崩れたり組み替わったりする印象を強調した．

特に「ザビ」では，グリッチ表現の時間変化を端末の動作と対応させた．端末の動きが小さい場合にはグリッチ用の時間値の進行を抑え，端末を動かした場合に時間変化を進める．これにより，静止している間に模様が自動的に変化し続ける状態を抑え，端末操作と表示異常との関係を知覚しやすくした．

一方，「ゆら」では，細かなノイズの密度や急激な破損を抑え，比較的大きな模様を用いる．また，前景側ほど大きく，背景側ほど小さく移動させることで，端末の傾きに応じたパララックスを強調する．これにより，端末を傾けた際に複数のグリッチレイヤが緩やかに移動する表現とした．

旧版では，各プリセットの違いを主としてRGBずれ量，ノイズサイズ，パララックス感度などの設定値によって生成していた．改良版では，これらに加えて，プリセットごとの描画特性を\texttt{NoiseProfile}として定義した．\texttt{NoiseProfile}では，細かなノイズの密度や大きさ，粗い矩形の密度・透明度，ちらつき，帯状ノイズの強度，RGBずれ，パララックスなどを個別に設定する．

これにより，単純な設定値の大小だけでなく，模様の粒度，密度，崩れ方，端末操作への反応まで含めて，「ザビ」と「ゆら」の特徴を区別できるようにした．

また，配色についても改良を行った．モノクロ系プリセットでは，基本となるパレットのRGB値から輝度を算出し，各色成分を同じ値へ変換することで，基本的な模様を無彩色化した．カラー系では，比較的鮮やかなパレットを使用する．ただし，モノクロ系についても，後段で適用されるRGBずれなどによって部分的に色が現れる場合がある．

描画レイヤの構成についても変更した．旧版では基本となる5レイヤに複数の補助レイヤを加えていたが，改良版では補助レイヤを整理し，「ザビ」にのみ前景の破片レイヤを追加する構成とした．本実装で使用する5基本レイヤの構成では，1フレームあたりの描画呼び出し回数は，旧版の10回に対して，「ゆら」では8回，「ザビ」では9回となる．

ただし，描画回数の削減だけで処理負荷が一定割合低下したことを意味するものではない．改良版では新たな破損表現に伴う計算も追加されているため，実際の消費電力や処理時間の改善量については別途測定が必要である．


\section{UIの改良}

予備評価時のシステムでは，8種類のプリセットを番号によって選択するUIに加え，RGBずれ量，ノイズサイズ，パララックス感度などを個別に変更するための複数の調整項目を設けていた．

この構成では，利用者が細かな表現を変更できる一方，操作項目が多く，現在どのような設定を使用しているのかを把握しにくいという問題があった．また，評価において複数の条件を比較する場合，各利用者が設定値を自由に変更すると提示条件を統一しにくい．

そこで改良版の評価用システムでは，個別のパラメータ調整UIを削除し，「ザビモノ」，「ザビカラー」，「ゆらモノ」，「ゆらカラー」の4種類を2×2に配置した選択画面へ変更した．利用者はプリセットを選択した後，ライブ壁紙として設定するだけで利用できる．

各プリセットには，RGBずれ量，ノイズサイズ，ちらつき，帯状ノイズの幅，パララックス感度などの値をあらかじめ設定している．これにより，利用者が多数の数値を調整することなく，特徴の異なるグリッチ表現を選択できるようにした．

一方，この変更によって，旧版と比較して利用者が個々のパラメータを細かく調整する自由度は低下している．改良版では，まず異なる特徴を持つ複数の表現から選択する方式に限定し，選好の違いを評価しやすい構成を優先した．


\section{プリセット選択によるパーソナライズ}

本研究では，予備評価の結果から，利用者によって好まれるグリッチ表現が異なる可能性に着目した．そこで改良版では，一つの表現を全利用者へ共通して提示するのではなく，複数の表現から利用者自身が選択できる構成とした．

現行の評価用システムにおけるパーソナライズは，グリッチを構成する各パラメータを利用者が自由に設定する方式ではなく，あらかじめ特徴の異なる4種類のプリセットから好みの表現を選択する方式である．

4種類は，「ザビ／ゆら」という端末操作への反応方式と，「モノ／カラー」という配色の組合せによって構成する．これにより，例えば，強い破損表現を好む利用者は「ザビ」を，緩やかな移動を好む利用者は「ゆら」を選択でき，さらに配色についてもモノクロ系とカラー系から選択できる．

この方式は，細かなパラメータをすべて利用者へ提示する場合と比較して自由度は低いものの，表現の違いを理解しやすく，少ない操作で自身の好みに近い表現を選択できるという利点を持つ．

なお，初期版ではパラメータを個別に変更する機能も実装していた．そのため，将来的には，プリセットによる大まかな選択を行った後，必要に応じて各パラメータを微調整する方式へ拡張することも可能である．


\section{評価方法の再設計}

予備評価では，8種類のプリセットを6名の参加者に固定順で提示し，各壁紙について「好みではない」から「好み」までの5段階評価を求めた．この結果から参加者間で選好が異なる傾向を確認できた一方，いくつかの課題も明らかとなった．

第一に，参加者数が6名と少なく，同一ゼミに所属する近い年齢層の学生に限定されていた．そのため，得られた結果を一般的な利用者全体の傾向として扱うことはできない．

第二に，8種類のプリセットでは複数の設定値を同時に変更していたため，どの視覚的特徴が評価へ影響したのかを明確に分離することが困難であった．例えば，あるプリセットが高く評価された場合でも，その理由が配色，ノイズサイズ，パララックス，ちらつきなどのどの要素にあるのかを判断できなかった．

第三に，評価項目が5段階の好み評価のみであり，参加者が評価した理由を取得していなかった．そのため，評価値の違いは確認できても，利用者がどのような特徴を好ましい，あるいは好ましくないと感じたのかを分析することはできなかった．

また，刺激は全参加者に対して同じ順序で提示しており，著者が端末を操作して参加者へ提示する形式であった．そのため，提示順による影響や，参加者自身が端末を操作した場合の印象についても十分に検討できていない．

これらの課題を踏まえ，次の評価では，改良した4種類のプリセットを用いて，利用者ごとの選好の違いを改めて確認するとともに，どのような視覚的特徴や動きが評価に関係するのかをより詳細に検討する．